\section{Current Limits and Failure Behavior}\label{sec:limits}

The most important current limit is the event cap in the lowerer.
\texttt{MAX\_EVENTS} is 20,000.
A program that generates 20,001 note events fails during lowering with the diagnostic
\texttt{program exceeds 20000 event limit (has 20001)}.
This is a deliberate compile-time safety boundary.
The cap provides a termination guarantee for programs whose iteration count is not statically bounded: Motivo accepts
\texttt{for (;;)} and other loops whose trip count semantic analysis cannot prove finite, so lowering must stop after a
fixed number of emitted events rather than unrolling indefinitely.
The value 20{,}000 is an implementation limit that stops unrolling with a diagnostic once reached; it is not
a claim about
maximum composition length in principle.

\subsection{Data layout and retained state}\label{detail-subsec:memory-footprint}

The event cap is the primary resource safeguard; memory use otherwise scales with the materialized timeline rather than
with abstract control flow.
The AST stores statements and expressions in \texttt{std::variant} sum types with \texttt{std::visit}
dispatch; recursive
structure uses \texttt{std::unique\_ptr} nodes while symbol and annotation tables live in contiguous vectors indexed by
opaque IDs.
The lowering IR holds only atomic fields per event (pitch, start beat, duration, velocity) in tightly packed
\texttt{struct}s.
This layout keeps traversal over contiguous data and limits retained state to the parsed tree plus the expanded event
list bounded by \texttt{MAX\_EVENTS}, rather than interpreter frames or runtime scheduling state.
